\scanpage{142}\mspage{72}

pour son propre mérite. La question bien sûr ne se pose que $Z'_n\ne\{1\}$, i.e. si $n\notin\{2,3,4,6\}$, le premier cas est donc celui où
\[
n=5,\qquad \Lambda'_{n}\simeq\deleted{\unsure{\Lambda}}\mathbb F_5^*/\pm1\simeq\mathbb Z/2\mathbb Z,\tag*{(37)}
\]
considérons une structure de 2-polyèdre connexe triangulé ($=$ faces d'ordre 3), dont les sommets sont \inserted{tous} d'ordre 5. On sait qu'il n'y a que l'icosaèdre (structure orientable à 12 sommets) et l'icosaèdre gauche (structure non orientée). L'opération envisagée ici (pour $\unsure{\nu|v}$ l'élément non trivial de $\Lambda'_5$) est celle qui fait passer d'un icosaèdre à l'icosaèdre associé -- l'ens $\Pi$ (dans le cas orienté) étant la structure de \uline{bi-icosaèdre}.

L'icosaèdre peut être décrit à l'aide d'un vectoriel $M$ de rang 2 sur $\mathbb F_5$, comme $(S_M,F_M)$, et satisfait donc à toute l'axiomatique du présent §. Mais l'icosaèdre gauche ? Notons que le commutant de $G_0$ dans ce cas est réduit à l'identité, donc $Z=1$, donc on ne s'attendra pas à ce que $Z\longrightarrow\deleted{\unsure{Z}}\Lambda_n^{'\!*}$ soit surjectif -- pourtant l'opération $A_0\longrightarrow A_0^{(\nu)}$ a un sens dans ce cas là ! Cela

% archive sheet number: 141
